\section{Introduction}

Microsimulation models of tax and transfer policy require individual records
that span domains no single survey covers
\citep{bourguignon2006microsimulation, sutherland2013euromod}. The Current
Population Survey observes employment and income but not wealth; the Survey of
Consumer Finances observes wealth in detail but lacks the \cps{}'s coverage
and program detail; administrative tax records observe income components
precisely but only for filers, and with narrow demographics. Model builders
bridge these gaps with imputation: fit a conditional model of the missing
variables on a donor survey, then draw values for the receiver's records
\citep{little2002statistical, dorazio2006statistical}.

Two properties of this setting distinguish it from generic missing-data
problems. First, both files are \emph{weighted samples}: each record stands in
for a different number of population units, and the quantity of interest is a
population distribution, not a sample one. An imputation model fit to
unweighted records recovers the sample conditional, which differs from the
population conditional exactly when the design (or a prior calibration)
correlates weights with outcomes. Second, the imputed file is not an endpoint.
It is \emph{reweighted downstream} --- calibrated to administrative margins
\citep{deville1992calibration, hainmueller2012entropy}, filtered to
subpopulations, stressed under reforms --- so an imputation that looks
acceptable under the weights it shipped with can fail under weights it never
saw. A rare donor record carrying an extreme value at near-zero weight is
invisible in the shipped aggregates; a later reweighting that assigns it mass
turns it into a distortion of the population total. We refer to such records
as landmines, and we measure exposure to them directly with a reweighting
stress test (\todo{fragility results reference}).

This paper presents and evaluates the imputation operator of \populace{},
\policyengine{}'s open-source microdata stack \citep{populace2026}. The
estimator combines three design choices, each motivated by one of the failure
modes above: a \emph{weighted bootstrap} that materializes survey weights into
the training data of each forest, so the learned conditionals are population
conditionals; \emph{regime gates} that classify each record's sign class
(negative, zero, positive) before magnitudes are modeled, so zero-inflated and
sign-mixed variables --- ubiquitous in economic microdata
\citep{mullahy1986specification, lambert1992zero} --- retain their mass
structure; and \emph{sequential chaining} across imputed variables, so joint
structure among targets survives the imputation
\citep{raghunathan2003multiple, vanbuuren2011mice}. The interface enforces the
first choice: fits are weighted by construction, and an unweighted fit must be
requested explicitly.

We make three contributions:

\begin{enumerate}
\item \textbf{A weight-aware imputation estimator, available standalone.}
We describe the estimator and its interface contract, and release it as
\populacefit{}, installable and usable on plain data frames independently of
the rest of the \populace{} stack.
\item \textbf{A benchmark against standard practice with attribution.} We
compare the estimator against unweighted and weighted quantile regression
forests \citep{meinshausen2006quantile}, ordinary least squares imputation
\citep{vonhippel2007should}, linear quantile regression
\citep{koenker1978regression}, and nearest-neighbour hot-deck statistical
matching \citep{andridge2010review, dorazio2021statistical}, under a paired,
repeated holdout protocol with weighted distributional metrics. Ablations
knock out one design choice at a time, attributing performance to weighting,
gating, and chaining separately.
\item \textbf{A reweighting stress test.} We introduce a fragility diagnostic
that measures the worst-case single-record contribution to a population
aggregate under admissible reweightings of the imputed file, formalizing the
landmine failure mode; we evaluate every method under it.
\end{enumerate}

The policy stakes are concrete. \todo{One-paragraph SSI Savings Penalty
Elimination Act exhibit: without imputed wealth the baseline caseload and the
reform are unsimulable; with weight-blind imputation the caseload is
misestimated; numbers from the task-1 sweep.}

The rest of the paper proceeds as follows. Section~2 situates the estimator in
the survey-imputation literature. Section~3 specifies the estimator, the
ablation design, the metrics, and the protocol. Section~4 describes the four
benchmark tasks and their data. Section~5 reports results, Section~6
discusses limitations, and Section~7 concludes.
